\documentclass[10pt,a4paper]{amsart}
\usepackage[margin=19mm]{geometry}
\usepackage{amsmath,amssymb,mathtools}
\usepackage[scr=rsfs,cal=euler]{mathalfa}
\usepackage{xfrac}
\usepackage[unicode,hidelinks,bookmarks=false]{hyperref}
\newcommand{\ie}{i.e.\ }
\newcommand{\cf}{cf.\ }
\newcommand{\resp}{respectively\ }
\newcommand{\Subsection}{\subsection}
\providecommand{\nobreakdash}{\nobreak\hbox{-}}
\setlength{\emergencystretch}{2em}
\multlinegap0pt
\usepackage{bm}
\usepackage{bbm}
\usepackage{dsfont}
\usepackage{xspace}
\usepackage{cancel}
\usepackage{mathtools}
\usepackage{amsthm}
\makeatletter
\@ifundefined{theorem}{\newtheorem{theorem}{Theorem}}{}
\@ifundefined{corollary}{\newtheorem{corollary}[theorem]{Corollary}}{}
\@ifundefined{definition}{\newtheorem{definition}[theorem]{Definition}}{}
\@ifundefined{deflem}{\newtheorem{deflem}[theorem]{Definition-Lemma}}{}
\@ifundefined{lemma}{\newtheorem{lemma}[theorem]{Lemma}}{}
\@ifundefined{proposition}{\newtheorem{proposition}[theorem]{Proposition}}{}
\@ifundefined{remark}{\newtheorem{remark}[theorem]{Remark}}{}
\makeatother
\DeclareMathOperator{\ores}{Ores}
\DeclareMathOperator{\pano}{PanOres}
\newcommand{\Ca}{\mathcal{C}} %field of constants in algebraic case
\newcommand{\Cq}{\mathcal{C}} %field of constants in Tate case
\newcommand{\Z}{\mathbb{Z}} % integers
\newcommand{\bk}{\mathcal{K}} %base field notation -- mnemonic: "base k"
\newcommand{\ec}{\mathcal{E}} % elliptic curve
\newcommand{\idec}{{o}} % identity point of elliptic curve
\newcommand{\orbita}{{\ec/\tau}} %orbits in algebraic case
\newcommand{\za}[2]{\zeta^{\mathcal{U}}_{( {#1} , {#2} )}}
\title[Panorbital residues and elliptic summability]{Panorbital residues and elliptic summability}
\author{Carlos E. Arreche, Matthew W. Babbitt}
\date{Source: arXiv:2508.18247v1 (CC BY 4.0)}
\begin{document}
\maketitle

\noindent\emph{Source and licence.} Panorbital residues and elliptic summability. Version arXiv:2508.18247v1 (CC BY 4.0), distributed under the Creative Commons Attribution 4.0 International licence, as recorded in the arXiv licence metadata for this exact version and shown on the abstract page https://arxiv.org/abs/2508.18247v1. Full rights notice: licenses/arxiv-2508.18247-CC-BY-4.0.txt.\par\smallskip

\noindent\textit{Editorial scope.} Counted unit: the Lemma whose source label is \texttt{lem:pano-pinning-au} in the article (source0.tex lines 762--770, source label \texttt{lem:pano-pinning-au}), delivered together with the article's own complete original proof of it (source0.tex lines 772--812, one \texttt{proof} environment of 41 lines). No mathematics is written, repaired or re-derived: statement and proof are the article's own text line for line. Required context retained uncounted: eq:uniformizers (lines 507--507); def:uniformizers (lines 514--516); def:oresa (lines 518--519); eq:u-coeff-effect (lines 521--526); rem:ores-uniformizers (lines 521--526); deflem:a-zeta-difference (lines 544--550); def:a-zeta-differences (lines 556--563); def:pinninga (lines 570--572); deflem:phia (lines 576--590); eq:d-r-independence (lines 604--605); lem:d-r-independence (lines 604--605); def:structural-constants (lines 631--640); lem:zetaExp-a (lines 666--671); cor:zetaExp-a (lines 679--680); prop:diff-pano (lines 685--688). Every definition, display and label of those lines is retained unchanged. Omitted from this package and not redistributed: 1--506, 508--513, 517--517, 520--520, 527--543, 551--555, 564--569, 573--575, 591--603, 606--630, 641--665, 672--678, 681--684, 689--761, 771--771, 813--1257. Nothing inside a retained excerpt is omitted or altered: every retained body is delivered exactly as the article's lines are written, so no omission receipt of any kind accompanies this package. Citations: the retained excerpts carry no citation command at all, so no bibliography entry of the article is transcribed and no reference is redistributed. Reference adaptation: none was needed and none was made; every reference inside the delivered unit resolves inside it, so no unresolved reference or citation is produced. Printed slips kept literal: no printed word, symbol, hypothesis or number of the counted statement or of its proof was altered. Layout and class: the article's own class is \texttt{elsarticle} with packages including amsmath, amssymb, amsthm, enumerate, microtype, xcolor, marginnote, bm, todonotes, hyperref, ifpdf; the wrapper is the lane's pinned \texttt{amsart} editorial wrapper at 10pt on A4, which restates the theorem-like environments the retained text needs and adds the article's own macro definitions that the retained text needs (\texttt{\textbackslash Ca}, \texttt{\textbackslash Cq}, \texttt{\textbackslash Z}, \texttt{\textbackslash bk}, \texttt{\textbackslash ec}, \texttt{\textbackslash idec}, \texttt{\textbackslash orbita}, \texttt{\textbackslash ores}, \texttt{\textbackslash pano}, \texttt{\textbackslash za}), with extra wrapper packages (bm, bbm, dsfont, xspace, cancel, mathtools). The delivered numbering is the unit's own. Assets: no retained span contains a figure environment, a graphics-inclusion command, a PDF-page inclusion, a table or a TikZ picture; no image file is copied into this package, the package declares no asset dependency and no image file is redistributed with it. Licence: version arXiv:2508.18247v1 of this article is distributed under the Creative Commons Attribution 4.0 International licence, as recorded in the arXiv licence metadata for this exact version and shown on the abstract page https://arxiv.org/abs/2508.18247v1. A full rights notice is delivered at licenses/arxiv-2508.18247-CC-BY-4.0.txt. Characters: every retained excerpt is pure ASCII, so the delivered engine, XeLaTeX with its default Latin Modern text font, reports no missing character.
Relative to a given choice of local uniformizers \begin{equation}\label{eq:uniformizers}\mathcal{U}=\bigl\{u_\alpha \in\bk \ \big| \ \nu_\alpha(u_\alpha)=1 \ \text{for every} \ \alpha\in\ec(\Ca)\bigr\},\end{equation} for any given $f\in\bk$ we write $c_k^\mathcal{U}(f,\alpha)\in\Ca$ for the coefficients occurring in the non-positive part of the formal Laurent series expansion of $f$ in $\Ca((u_\alpha))$:

\begin{definition}[$\tau$-compatible uniformizers] \label{def:uniformizers}
    A set $\mathcal{U}=\{u_\alpha\ | \ \alpha\in\ec(\Ca)\}$ of local uniformizers as in \eqref{eq:uniformizers} is $\tau$-\emph{compatible} if $\tau(u_{\alpha\oplus s})=u_\alpha$ for every $\alpha\in\ec(\Ca)$.
\end{definition}

\begin{definition}[Orbital residues - algebraic version] \label{def:oresa} For a given $\tau$-compatible set of local uniformizers $\mathcal{U}$, the \emph{orbital residue} of $f\in\bk$ at the orbit $\omega\in\orbita$ of order $k\in\mathbb{N}$ relative to $\mathcal{U}$ is \[\ores^\mathcal{U}(f,\omega,k):=\sum_{\alpha\in\omega}c_k^\mathcal{U}(f,\alpha).\]  
\end{definition}

\begin{remark}\label{rem:ores-uniformizers} In principle, it is a relatively simple matter to compute the effect of the choice of $\tau$-compatible set of local uniformizers $\mathcal{U}$ on the orbital residues in Definition~\ref{def:oresa}. Indeed, for another given choice $\mathcal{U}'=\{u_\alpha' \ | \ \alpha\in\ec(\Ca)\}$ of $\tau$-compatible local uniformizers, one can compute explicitly (as many terms as desired of) the unique formal power series \[h_\omega(u)=\sum_{j\geq 0}h_{\omega,j}u^j\in\Ca[[u]]\] with $h_{\omega,0}\neq 0$ such that $u_\alpha=u_\alpha'/h_\omega(u_\alpha')\in\Ca[[u_\alpha']]$ for every $\alpha\in\omega$ simultaneously. Writing $h_{\omega,j}^{(m)}$ for the coefficients occurring in $\bigl(h_{\omega}(u)\bigr)^m=\sum_{j\geq 0}h_{\omega,j}^{(m)}u^j$, a straightforward computation yields that, for each $f\in\bk$, \begin{gather}\label{eq:u-coeff-effect}
    c_k^{\mathcal{U}'}(f,\alpha)=\sum_{m\geq k}h_{\omega,m-k}^{(m)}\cdot c_m^\mathcal{U}(f,\alpha)\qquad\text{for every} \quad\alpha\in\omega\quad \text{and} \quad k\geq 0,
    \intertext{and therefore for each $\omega\in\orbita$ and $k\geq 1$ we have}
    \label{eq:u-ores-effect} \ores^{\mathcal{U}'}(f,\omega,k)=\sum_{m\geq k}h_{\omega,m-k}^{(m)}\cdot\ores^\mathcal{U}(f,\omega,m).
\end{gather}
\end{remark}

\begin{deflem}\label{deflem:a-zeta-difference} For $\mathcal{U}$ a $\tau$-compatible set of local uniformizers as in Definition~\ref{def:uniformizers}, there exists a unique $\za{1}{0} \in\mathcal{L}\bigl([s]+[\idec]\bigr)$ such that
\begin{enumerate}
    \item $c_1^\mathcal{U}\bigl(\za{1}{0},s\bigr)=1$;
    \item $c_1^\mathcal{U}\bigl(\za{1}{0},\idec\bigr)=-1$; and
    \item $c_0^\mathcal{U}\bigl(\za{1}{0},\idec\bigr)=0$.
\end{enumerate}
\end{deflem}

\begin{definition}\label{def:a-zeta-differences}
    With $\za{1}{0}\in\bk$ as in Definition~\ref{deflem:a-zeta-difference}, for $\ell,m\in\Z$ we define the elements $\za{\ell}{m}\in\bk$ as follows.
    \[\za{\ell}{m}:=\begin{cases}
        \displaystyle \sum_{i=m}^{\ell-1} \tau^{-i}\bigl(\za{1}{0}\bigr) & \text{if} \ \ell > m; \\
        0 \vphantom{\dfrac{A}{A}}& \text{if} \ \ell=m;\\
        \displaystyle -\sum_{i=\ell}^{m-1} \tau^{-i}\bigl(\za{1}{0}\bigr) & \text{if} \ \ell<m.     
    \end{cases}\]
\end{definition}

\begin{definition}[$\tau$-pinnings - algebraic version]\label{def:pinninga}
    A $\tau$-\emph{pinning} of $\ec$ is a choice $\mathbf{P}=(\mathcal{R},\mathcal{U})$, consisting of a set $\mathcal{R}=\{\alpha_\omega\in\ec(\Cq) \ | \ \omega\in\orbita\}$ of representatives $\alpha_\omega\in\omega$ for each orbit $\omega\in\orbita$, and a set of $\tau$-compatible local uniformizers $\mathcal{U}=\{u_\alpha \ | \ \alpha\in\ec(\Ca)\}$ as in Definition~\ref{def:uniformizers}.
\end{definition}

\begin{deflem} \label{deflem:phia}Let $\mathbf{P}=(\mathcal{R},\mathcal{U})$ be any $\tau$-pinning of $\ec$ as in Definition~\ref{def:pinninga}, with $\mathcal{R}=\{\alpha_\omega\}_{ \omega\in\orbita}$. 
\begin{enumerate}
    
    \item Given $\omega\in\orbita$ such that $\alpha_\omega\neq\idec$, $n\in\Z$, and $k\geq 1$, there exists a unique $\varphi^\mathbf{P}_{\omega,n,k}\in\mathcal{L}\bigl(k[\alpha_\omega\oplus ns]+[ns]\bigr)$ such that \begin{enumerate}
    \item $c_k^\mathcal{U}\bigl(\varphi^\mathbf{P}_{\omega,n,k},\alpha_\omega\oplus ns\bigr)=1$; 
    \item $c_j^\mathcal{U}\bigl(\varphi^\mathbf{P}_{\omega,n,k},\alpha_\omega\oplus ns\bigr)=0$ for $j\in\mathbb{N}-\{k\}$; and 
    \item $c_0^\mathcal{U}\bigl(\varphi^\mathbf{P}_{\omega,n,k},ns\bigr)=0$.
    \end{enumerate}
    \item Given $n\in\Z$ and $k\geq 2$, there exists a unique $\hat{\varphi}^\mathcal{U}_{\Z s,n,k}\in\mathcal{L}\bigl(k[ns]\bigr)$ such that \begin{enumerate}
    \item $c_k^\mathcal{U}\bigl(\hat{\varphi}^\mathcal{U}_{\Z s,n,k},ns\bigr)=1$; and
        \item $c_j^\mathcal{U}\bigl(\hat{\varphi}^\mathcal{U}_{\Z s,n,k},ns\bigr)=0$ for every $j\in\Z_{\geq 0}-\{1,k\}$.
    \end{enumerate}
\end{enumerate}
In case $\alpha_{\Z s}=\idec$, we write $\varphi^\mathbf{P}_{\Z s,n,1}:=0$ and $\varphi^{\mathbf{P}}_{\Z s,n,k}:=\hat{\varphi}^{\mathcal{U}}_{\Z s,n,k}$ for every $n\in\Z$ and $k\geq 2$.
\end{deflem}

\begin{lemma}\label{lem:d-r-independence} Let $\mathcal{U}=\{u_\alpha\}_{\alpha\in\ec(\Ca)}$ be a $\tau$-compatible set of local uniformizers and suppose that $\varphi\in\mathcal{L}\bigl(k[\alpha]+[ns]\bigr)$ for some $n\in\Z$, $k\geq 1$, and $ns\neq\alpha\in\ec(\Ca)$ such that \begin{equation}\label{eq:d-r-independence}c_j^\mathcal{U}(\varphi,\alpha)=\begin{cases} 1 & \text{for} \ j=k; \ \text{and} \\0 & \text{for} \ j\in\mathbb{N}-\{ k\}.\end{cases}\end{equation} Then $d^\mathcal{U}_{\omega,k}:=-c_1^\mathcal{U}(\varphi,ns)$ depends only on the orbit $\omega:=\alpha\oplus \Z s$ and the degree~$k$, and not on the particular choices of $\alpha\in\omega$, $n\in\Z$, nor $\varphi\in\mathcal{L}\bigl(k[\alpha]+[ns]\bigr)$ satisfying \eqref{eq:d-r-independence}. Moreover, $d^\mathcal{U}_{\Z s,1}=1$, and $c_1^\mathcal{U}\bigl(\hat{\varphi}^\mathcal{U}_{\Z s,n,k},ns\bigr)=-d^\mathcal{U}_{\Z s, k}$ for every $n\in\Z$ and $k\geq 2$, where $\hat{\varphi}^\mathcal{U}_{\Z s, n, k}$ is given as in Definition~\ref{deflem:phia}(2).
\end{lemma}

\begin{definition}[Structural constants of a $\tau$-pinning] \label{def:structural-constants}
    With notation as in Definition~\ref{def:a-zeta-differences} and Definition~\ref{deflem:phia}, we define the following structural constants arising from the $\tau$-pinning $\mathbf{P}=(\mathcal{R},\mathcal{U})$ of $\ec$.
     \begin{enumerate}
        \item For $\omega\in\orbita$ and $k\geq 1$, we set (for some/any $n\in\Z$) \[d_{\omega,k}^{\mathcal{U}}:=\begin{cases}
        -c_1^\mathcal{U}\bigl(\varphi^\mathbf{P}_{\omega,n,k},ns\bigr) & \text{if} \ \alpha_{\omega}\neq\idec \ \text{or} \ k\geq 2;\\
            1 & \text{if} \ \alpha_{\omega}=\idec \ \text{and} \ k=1.
        \end{cases}\]
        \item For $\omega\in\orbita$, $n\in \Z$, and $k\geq 1$, we set \[e^\mathbf{P}_{\omega,n,k}:=c_0^\mathcal{U}\Bigl(\varphi^\mathbf{P}_{\omega,n,k} + d^\mathcal{U}_{\omega,k}\cdot\za{n}{0},\ \idec\Bigr).\]
    \end{enumerate}
\end{definition}

\begin{lemma}\label{lem:zetaExp-a}
    Let $\mathbf{P}=(\mathcal{R},\mathcal{U})$ be any $\tau$-pinning of $\mathcal{E}$ and let $f\in\bk$. With notation as in Definition~\ref{def:a-zeta-differences},  Definition~\ref{deflem:phia}, and Definition~\ref{def:structural-constants},
    \begin{equation}\label{eq:zetaExp-a}
    f=\pano^\mathbf{P}(f,0)+\sum_{\omega\in\orbita}\sum_{n\in\Z}\sum_{k\geq 1}c^\mathcal{U}_k(f,\alpha_\omega\oplus ns)\cdot\Bigl(\varphi^\mathbf{P}_{\omega,n,k}+d^\mathcal{U}_{\omega,k}\cdot\za{n}{0}\Bigr).
    \end{equation}
\end{lemma}

\begin{corollary} \label{cor:zetaExp-a}For any $\tau$-compatible set of local uniformizers $\mathcal{U}$ and for any $f\in\bk$, \[\sum_{\omega\in\orbita}\sum_{k\geq 1}\ores^\mathcal{U}(f,\omega,k)\cdot d^\mathcal{U}_{\omega,k}=0.\]
\end{corollary}

\begin{proposition} \label{prop:diff-pano}
    Let $\mathbf{P}=(\mathcal{R},\mathcal{U})$ be a $\tau$-pinning of $\ec$ with $\mathcal{U}=\{u_\alpha\}_{\alpha\in\ec(\Ca)}$ as in Definition~\ref{def:uniformizers}. For the unique invariant differential $0\neq\varpi\in H^0(\ec,\Omega^1_{\ec/\Ca})$ such that the valuation $\nu_\idec(du_\idec-\varpi)\geq 1$, we have
    \[\pano^\mathbf{P}(f,1)=\sum_{\omega\in\orbita}\sum_{n\in\Z}n\cdot\mathrm{res}(f\varpi,\alpha_\omega\oplus ns).\]
\end{proposition}

\begin{lemma}\label{lem:pano-pinning-au} Let $\mathbf{P}=(\mathcal{R},\mathcal{U})$ be a $\tau$-pinning of $\ec$ as in Definition~\ref{def:pinninga}, with $\mathcal{R}=\{\alpha_\omega\}_{\omega\in\orbita}$ and $\mathcal{U}=\{u_\alpha\}_{\alpha\in\ec(\Ca)}$. Let $\mathcal{U}'=\{u_\alpha'\}_{\alpha\in\ec(\Ca)}$ be another choice of $\tau$-compatible set of local uniformizers, and let $h_{\Z s,0},h_{\Z s,1}\in\Ca$ be given by $h_{\Z s,\ell}:=c^{\mathcal{U}'}_{1-\ell}\bigl(u_\idec^{-1},\idec\bigr)$ for $\ell\in\{0,1\}$. Then for $\mathbf{P}':=(\mathcal{R},\mathcal{U}')$ we have
\begin{equation}\label{eq:pano-pinning-au1}
    \pano^{\mathbf{P}'}(f,1) =h_{\Z s,0}\cdot \pano^\mathbf{P}(f,1).\end{equation}Moreover, provided that $\alpha_{\Z s}\neq\idec$, we also have
    \begin{equation}
    \pano^{\mathbf{P}'}(f,0) =\pano^\mathbf{P}(f,0)-h_{\Z s,1}\cdot\pano^\mathbf{P}(f,1),\label{eq:pano-pinning-au0-nice}\end{equation}whereas in the exceptional situation where $\alpha_{\Z s}=\idec$ we have
    \begin{multline}\label{eq:pano-pinning-au0-mean}
    \pano^{\mathbf{P}'}(f,0) =\pano^\mathbf{P}(f,0)-h_{\Z s,1}\cdot\pano^\mathbf{P}(f,1)\\ +\sum_{n\in\Z}\Bigl(c^{\mathcal{U}'}_0(f,ns)-c^{\mathcal{U}}_0(f,ns)\Bigr).
\end{multline}
\end{lemma}

\begin{proof} Recall that one can compute $h^{(m)}_{\omega,j}\in\Ca$ for $m,j\geq 0$ and $\omega\in\orbita$ such that each $c_k^{\mathcal{U}'}(\varphi,\alpha)=\sum_{m\geq k}h^{(m)}_{\omega,m-k}\cdot c_m^{\mathcal{U}}(\varphi,\alpha)$ for every $\varphi\in\bk$, $\alpha\in \omega$, and $k\geq 1$ simultaneously as in \eqref{eq:u-coeff-effect} in~Remark~\ref{rem:ores-uniformizers}, where we denoted $h^{(1)}_{\omega,k}=h_{\omega,k}$ consistently with the usage of $h_{\Z s,0}$ and $h_{\Z s,1}$ in the statement of Lemma~\ref{lem:pano-pinning-au}.

Letting $\varpi\in H^0(\ec,\Omega^1_{\ec/\Ca})$ such that the valuation $\nu_\idec(du_\idec-\varpi)\geq 1$ as in Proposition~\ref{prop:diff-pano}, we find that $\varpi':=h_{\Z s,0}\cdot\varpi$ similarly satisfies $\nu_\idec(du_\idec'-\varpi')\geq 1$ (cf.~Remark~\ref{rem:ores-uniformizers}). We obtain \eqref{eq:pano-pinning-au1} by a twofold application of Proposition~\ref{prop:diff-pano}:
\begin{multline*}\pano^{\mathbf{P}'}(f,1)=\sum_{\omega\in\orbita}\sum_{n\in\Z}n\cdot\mathrm{res}(f\varpi',\alpha_\omega\oplus ns)\\=h_{\Z s,0}\sum_{\omega\in\orbita}\sum_{n\in\Z}n\cdot\mathrm{res}(f\varpi,\alpha_\omega\oplus ns)=h_{\Z s,0}\pano^\mathbf{P}(f,1).\end{multline*}
    
 By Lemma~\ref{lem:zetaExp-a}, \eqref{eq:pano-pinning-au0-nice} will follow from the forthcoming computation, valid in case $\alpha_{\Z s}\neq \idec$, that \begin{gather}\label{eq:pano-pinning-au0-proof} \sum_{\omega\in\orbita}\sum_{n\in\Z}\sum_{k\geq 1}c^{\mathcal{U}'}_k(f,\alpha_\omega\oplus ns)\cdot\Bigl(\varphi^{\mathbf{P}'}_{\omega,n,k}+d^{\mathcal{U}'}_{\omega,k}\cdot\zeta^{\mathcal{U}'}_{(n,0)}\Bigr)\\ =h_{\Z s,1}\cdot\pano^\mathbf{P}(f,1)+\sum_{\omega\in\orbita}\sum_{n\in\Z}\sum_{k\geq 1}c^{\mathcal{U}}_k(f,\alpha_\omega\oplus ns)\cdot\Bigl(\varphi^{\mathbf{P}}_{\omega,n,k}+d^{\mathcal{U}}_{\omega,k}\cdot\zeta^{\mathcal{U}}_{(n,0)}\Bigr),\notag\end{gather} where $\varphi_{\omega,n,k}^{\mathbf{P}}$ and $\varphi_{\omega,n,k}^{\mathbf{P}'}$ are as in Definition~\ref{deflem:phia}, $\za{n}{0}$ and $\zeta^{\mathcal{U}'}_{(n,0)}$ are as in Definition~\ref{def:a-zeta-differences}, and $d^\mathcal{U}_{\omega,k}$ and $d^{\mathcal{U}'}_{\omega,k}$ are as in Definition~\ref{def:structural-constants}. To begin, \mbox{we find that}
 
 \begin{equation}\label{eq:zeta-u-effect}
        h_{\Z s,0}\cdot\zeta^{\mathcal{U}'}_{(n,0)}=\za{n}{0}+h_{\Z s,1}\cdot n;
    \end{equation}
    first in the basic case $n=1$ by comparing local expansions at $\Ca ((u_\idec'))$ and using the computations that $c_1^{\mathcal{U}'}(\za{1}{0},s)=h_{\Z s,0}$ and $c_0^{\mathcal{U}'}(\za{1}{0},\idec)=h_{\Z s,1}$ by \eqref{eq:u-coeff-effect} and Definition~\ref{deflem:a-zeta-difference}, and then for arbitrary $n\in\Z$ from the Definition~\ref{def:a-zeta-differences}. Similarly, comparing local expansions in $\Ca((u_{\alpha_\omega\oplus ns}'))$, we find for every $\omega\in\orbita$, $n\in\Z$, and $k\geq 1$, that
    \begin{equation}\label{eq:phi-u-effect}
        \sum_{j=1}^k h^{(k)}_{\omega,k-j}\cdot\varphi^{\mathbf{P}'}_{\omega,n,j}=\varphi^\mathbf{P}_{\omega,n,k}-c^{\mathcal{U}'}_0\bigl(\varphi^\mathbf{P}_{\omega,n,k},\,ns\bigr).
    \end{equation} Indeed, the difference between the two sides of \eqref{eq:phi-u-effect} belongs to $\mathcal{L}([ns])=\Ca$ by \eqref{eq:u-coeff-effect}, and the the condition that each $c^{\mathcal{U}'}_0(\varphi^{\mathbf{P}'}_{\omega,n,j},ns)=0$ in Definition~\ref{deflem:phia} yields the equality.
    Further comparing local expansions in $\Ca((u_{ns}'))$, we obtain from \eqref{eq:phi-u-effect} that \begin{equation}\label{eq:d-u-effect}
    \sum_{j=1}^k h^{(k)}_{\omega,k-j}\cdot d^{\mathcal{U}'}_{\omega,j}=h_{\Z s,0}\cdot d^\mathcal{U}_{\omega,k}.
    \end{equation} 


    Provided that we are in the non-exceptional case where $\alpha_\omega\neq\idec$, we compute from \eqref{eq:u-coeff-effect} that the constant correction term in \eqref{eq:phi-u-effect} is \begin{equation}\label{eq:phi-u-effect-constant}
        -c_0^{\mathcal{U}'}\bigl(\varphi^\mathbf{P}_{\omega,n,k},ns\bigr)=h_{\Z s,1}\cdot d^\mathcal{U}_{\omega,k}.
    \end{equation} Therefore, for each $\omega\in\orbita$ such that $\alpha_\omega\neq\idec$ and for each $n\in \Z$, we find that \begin{multline}\label{eq:pano-pinning-au0-long-computation-nice}
    \sum_{k\geq 1} c^{\mathcal{U}'}_k(f,\alpha_\omega\oplus ns)\cdot\Bigl(\varphi^{\mathbf{P}'}_{\omega,n,k}+d^{\mathcal{U}'}_{\omega,k}\cdot\zeta^{\mathcal{U}'}_{(n,0)}\Bigr)\\
      =\sum_{k\geq 1}\left(\sum_{m\geq k} h^{(m)}_{\omega,m-k}\cdot c^{\mathcal{U}}_m(f,\alpha_\omega\oplus ns)\right)\cdot\Bigl(\varphi^{\mathbf{P}'}_{\omega,n,k}+d^{\mathcal{U}'}_{\omega,k}\cdot\zeta^{\mathcal{U}'}_{(n,0)}\Bigr) \\
      =\sum_{k\geq 1}c_k^\mathcal{U}(f,\alpha_\omega\oplus ns)\cdot\sum_{j=1}^k h^{(k)}_{\omega,k-j}\cdot\Bigl(\varphi^{\mathbf{P}'}_{\omega,n,j}+d^{\mathcal{U}'}_{\omega,j}\cdot\zeta^{\mathcal{U}'}_{(n,0)}\Bigr)\\
      =\sum_{k\geq 1}c_k^\mathcal{U}(f,\alpha_\omega\oplus ns)\cdot\Bigl(\varphi^\mathbf{P}_{\omega,n,k}+h_{\Z s,1}d_{\omega,k}^{\mathcal{U}}+d^\mathcal{U}_{\omega,k}\za{n}{0}+d^{\mathcal{U}}_{\omega,k}h_{\Z s,1}n\Bigr),
    \end{multline} where the first equality follows from \eqref{eq:u-coeff-effect}, the second equality is a rearrangement, and the third equality is obtained from \eqref{eq:phi-u-effect}, \eqref{eq:phi-u-effect-constant},  \eqref{eq:d-u-effect}, and \eqref{eq:zeta-u-effect}. Thus, provided that $\alpha_{\Z s}\neq\idec$, summing over all $\omega\in\orbita$ and $n\in\Z$ yields \eqref{eq:pano-pinning-au0-proof}, in view of Corollary~\ref{cor:zetaExp-a}.

    It remains to prove \eqref{eq:pano-pinning-au0-mean} in the exceptional case where $\alpha_{\Z s}=\idec$. With $\hat{\varphi}^\mathcal{U}_{\Z s,n,k}=\varphi^\mathbf{P}_{\omega,n,k}$ and $\hat{\varphi}^{\mathcal{U}'}_{\omega,n,k}=\varphi^{\mathbf{P}'}_{\omega,n,k}$ as in Definition~\ref{deflem:phia}(2), we find from \eqref{eq:u-coeff-effect} that in this case the constant correction term in \eqref{eq:phi-u-effect} is \begin{equation}\label{eq:pano-pinning-au0-phihat-constant}
        -c_0^{\mathcal{U}'}\bigl(\hat{\varphi}^\mathcal{U}_{\Z s,n,k},ns\bigr)=h_{\Z s, 1}\cdot d^\mathcal{U}_{\Z s,k}-h^{(k)}_{\Z s,k},
    \end{equation} so that \eqref{eq:phi-u-effect}, which remains valid as stated in all cases and regardless of whether $\alpha_{\Z s}=\idec$, specializes in this case to
    \begin{equation}\label{eq:phihat-u-effect}
        \sum_{j=1}^k h^{(k)}_{\Z s,k-j}\cdot\hat{\varphi}^{\mathcal{U}'}_{\Z s,n,j}=\hat{\varphi}^\mathcal{U}_{\Z s,n,k}+h_{\Z s,1}\cdot d^\mathcal{U}_{\Z s,k}-h^{(k)}_{\Z s,k}.
    \end{equation} The relations \eqref{eq:zeta-u-effect} and \eqref{eq:d-u-effect} also remain valid, since they do not depend on the choice of $\mathcal{R}$ at all -- although one can also tortuously recover \eqref{eq:d-u-effect} for $\omega=\Z s$ by carefully comparing local $\Ca((u_{ns}'))$-expansions in \eqref{eq:phihat-u-effect} and using the fact that $d^{\mathcal{U}'}_{\Z s,1}=1=d^\mathcal{U}_{\Z s,1}$ by Lemma~\ref{lem:d-r-independence}. Thus the analogous computation to \eqref{eq:pano-pinning-au0-long-computation-nice} in this case yields that, for each $n\in\Z$,
\begin{gather}\label{eq:pano-pinning-au0-long-computation-mean} \sum_{k\geq 1}c^{\mathcal{U}'}_k(f,ns)\cdot\Bigl(\hat{\varphi}^{\mathcal{U}'}_{\Z s,n,k}+d^{\mathcal{U}'}_{\Z s,k}\cdot\zeta^{\mathcal{U}'}_{(n,0)}\Bigr)\\
=\sum_{k\geq 1}c^\mathcal{U}_k(f,ns)\cdot(\hat{\varphi}^\mathcal{U}_{\Z s,n,k}+d^\mathcal{U}_{\Z s,k}\cdot\za{n}{0}+h_{\Z s,1}\cdot(n+1)\cdot d^\mathcal{U}_{\Z s,k} -h^{(k)}_{\Z s,k}\Bigr).\notag
\end{gather} Thus in case $\alpha_{\Z s}=\idec$ we find that instead of \eqref{eq:pano-pinning-au0-long-computation-nice} we obtain
\begin{multline} \label{eq:pano-pinning-au0-long-computation-mean-result}
    \sum_{\omega\in\orbita}\sum_{n\in\Z}\sum_{k\geq 1}c^{\mathcal{U}'}_k(f,\alpha_\omega\oplus ns)\Bigl(\hat{\varphi}^{\mathcal{U}'}_{\omega,n,k}+d^{\mathcal{U}'}_{\omega,k}\cdot\zeta^{\mathcal{U}'}_{(n,0)}\Bigr)\\=h_{\Z s,1}\cdot\pano^\mathbf{P}(f,1)+\sum_{\omega\in\orbita}\sum_{n\in\Z}\sum_{k\geq 1}c^{\mathcal{U}}_k(f,\alpha_\omega\oplus ns)\cdot\Bigl(\hat{\varphi}^{\mathcal{U}'}_{\omega,n,k}+d^{\mathcal{U}}_{\omega,k}\cdot\zeta^{\mathcal{U}}_{(n,0)}\Bigr)\\-\sum_{n\in\Z}\sum_{k\geq 1}h^{(k)}_{\Z s,k}\cdot c^\mathcal{U}_k(f,ns).
\end{multline}The sought relation \eqref{eq:pano-pinning-au0-mean} follows from \eqref{eq:u-coeff-effect} because, for each $n\in\Z$, \[\sum_{k\geq 1}h^{(k)}_{\Z s,k}\cdot c^\mathcal{U}_k(f,ns)=c^{\mathcal{U}'}_0(f,ns)-c_0^\mathcal{U}(f,ns).\qedhere\]
\end{proof}

\end{document}
